\documentclass[11pt, a4paper]{article}
\usepackage[utf8]{inputenc}
\usepackage[T1]{fontenc}
\usepackage{amsmath}
\usepackage{amssymb}
\usepackage{amsfonts}
\usepackage{geometry}
\usepackage{enumitem}
\usepackage{fancyhdr}

% Page layout configuration
\geometry{left=1in, right=1in, top=1in, bottom=1in}
\setlength{\parindent}{0pt}
\setlength{\parskip}{1em}

% Header/Footer
\pagestyle{fancy}
\fancyhf{}
\lhead{\textbf{Partial Differential Equations (H)}}
\rhead{\textbf{2025 Fall Midterm}}
\cfoot{\thepage}

\begin{document}

\begin{center}
    \huge \textbf{PDE(H) Midterm}
    \vspace{0.5cm}

    \large 10:10 --- 12:10 $\quad$ Nov 6, 2025
\end{center}

\vspace{0.5cm}

\textbf{Question 1.}

Consider the initial value problem for the Burgers' equation:
\[
    u_t + u \cdot u_x = 0, \quad u(0, x) = \phi(x), \quad (t, x) \in [0, \infty) \times \mathbb{R}.
\]

\begin{enumerate}[label=(\alph*)]
    \item (10 points) Let
          \[
              \phi(x) = \begin{cases}
                  0,   & x \le 0, \\
                  x^2, & x > 0.
              \end{cases}
          \]
          Find a solution $u(t, x) \in C^{1,1}((0, \infty) \times \mathbb{R}) \cap C([0, \infty) \times \mathbb{R})$.

    \item (10 points) Let
          \[
              \phi(x) = \begin{cases}
                  1,       & x \le 0,     \\
                  1 - x^2, & 0 < x \le 1, \\
                  0,       & x > 1.
              \end{cases}
          \]
          Find the largest time $t_s$ such that all characteristics do not intersect.
\end{enumerate}
\textbf{[Total for Question 1: 20 points.]}

\noindent \rule[0.5pt]{\textwidth}{0.5pt}

\textbf{Question 2.}
\begin{enumerate}[label=(\alph*)]
    \item (5 points) Find all the eigen-pairs $(\lambda_n, X_n)$ of the Sturm-Liouville problem:
          \[
              X''(x) + \lambda X(x) = 0, \quad X'(0) = 0, \quad X(\pi) = 0.
          \]

    \item (10 points) Use Separation of Variables to solve the homogeneous heat equation:
          \[
              \begin{cases}
                  u_t - u_{xx} = 0,                                   & x \in (0, \pi), \ t > 0, \\
                  u(0, x) = 25 \cos \frac{x}{2} + \cos \frac{11x}{2}, & x \in [0, \pi],          \\
                  u_x(t, 0) = u(t, \pi) = 0,                          & t \ge 0.
              \end{cases}
          \]

    \item (5 points)  Use Duhamel's principle to solve the inhomogeneous heat equation:
          \[
              \begin{cases}
                  u_t - u_{xx} = \cos \frac{5x}{2},                   & x \in (0, \pi), \ t > 0, \\
                  u(0, x) = 25 \cos \frac{x}{2} + \cos \frac{11x}{2}, & x \in [0, \pi],          \\
                  u_x(t, 0) = u(t, \pi) = 0,                          & t \ge 0.
              \end{cases}
          \]

    \item (5 points) Suppose the initial condition is $$u(0, x) = \phi(x) \in C^1[0, \pi],$$ where $\phi'(0) = \phi(\pi) = 0$. Show that
          \[
              \lim_{t \to \infty} u(t, x) = v(x)
          \]
          for some function $v(x)$ independent of $\phi$, and compute the exact form of $v$.
\end{enumerate}
\textbf{[Total for Question 2: 25 points.]}

\noindent \rule[0.5pt]{\textwidth}{0.5pt}

\textbf{Question 3.}

Let $\Omega \subset \mathbb{R}^d$ be a bounded domain and $T > 0$. Let $b, c \in C(\bar{\Omega}_T)$ with $c \ge 0$.

\begin{enumerate}[label=(\alph*)]
    \item (10 points) Prove the following comparison principle: if $v \in C^{1,2}(\Omega_T) \cap C(\bar{\Omega}_T)$ satisfies
          \[
              \begin{cases}
                  v_t - \Delta v + b \cdot \nabla v + cv \ge 0, & (t, x) \in \Omega_T,            \\
                  v(0, x) \ge 0,                                & x \in \Omega,                   \\
                  v(t, x) \ge 0                                 & t > 0, \ x \in \partial \Omega,
              \end{cases}
          \]
          then $v(t, x) \ge 0$ on $\bar{\Omega}_T$.

    \item (10 points) Let $u \in C^{1,2}(\Omega_T) \cap C(\bar{\Omega}_T)$ solve
          \[
              \begin{cases}
                  u_t - \Delta u + b \cdot \nabla u + cu = f(t, x), & (t, x) \in \Omega_T,            \\
                  u(0, x) = \phi(x),                                & x \in \Omega,                   \\
                  u(t, x) = 0,                                      & t > 0, \ x \in \partial \Omega.
              \end{cases}
          \]
          Show that there exists a constant $C = C(T)$ such that
          \[
              \sup_{\bar{\Omega}_T} |u| \le C \left( \sup_{\Omega_T} |f| + \sup_{\Omega} |\phi| \right).
          \]
\end{enumerate}
\textbf{[Total for Question 3: 20 points.]}

\noindent \rule[0.5pt]{\textwidth}{0.5pt}

\textbf{Question 4.}

(15 points) For $\Omega = B_r(0) \subset \mathbb{R}^3$, let $u \in C^2(\Omega) \cap C(\bar{\Omega})$ solve
\[
    \begin{cases}
        -\Delta u = f, & x \in \Omega,          \\
        u = g,         & x \in \partial \Omega.
    \end{cases}
\]
where $f, g$ are continuous. Show that $u$ satisfies
\[
    u(0) = \frac{1}{4\pi r^2} \int_{\partial \Omega} g(x) \, dS(x) + \frac{1}{4\pi} \int_{\Omega} (|x|^{-1} - r^{-1}) f(x) \, dx.
\]
\textbf{[Total for Question 4: 15 points.]}

\noindent \rule[0.5pt]{\textwidth}{0.5pt}

\textbf{Question 5.}

(10 points) Let $u \in C^2(\Omega) \cap C^1(\bar{\Omega})$ solve
\[
    \begin{cases}
        -\Delta u + b(x) \cdot \nabla u + c(x)u = f(x), & \Omega,          \\
        \frac{\partial u}{\partial n} + \alpha(x)u = 0, & \partial \Omega,
    \end{cases}
\]
where $\alpha(x) \ge 0$, $c(x) \ge 1$, and $|b(x)| \le 1$. Show that there exists a constant $M$ such that
\[
    \int_{\Omega} |\nabla u(x)|^2 \, dx + \int_{\Omega} |u(x)|^2 \, dx \le M \int_{\Omega} |f(x)|^2 \, dx.
\]
Here, $|v|$ denotes the Euclidean norm of the vector $v \in \mathbb{R}^d$, i.e., $|v| := (v_1^2 + \cdots + v_d^2)^{1/2}$.

\textbf{[Total for Question 5: 10 points.]}

\noindent \rule[0.5pt]{\textwidth}{0.5pt}

\textbf{Question 6.}
\begin{enumerate}[label=(\alph*)]
    \item (3 points) Suppose that $u_* \in C^2(\bar{\Omega})$ minimizes the functional
          \[
              I[u] = \int_{\Omega} \sqrt{1 + |\nabla u|^2} \, dx, \quad \text{subject to } u|_{\partial \Omega} = g.
          \]
          Derive the equation that $u$ must satisfy (the Euler--Lagrange equation). \\
          \textit{Hint: Use the condition $I[u_* + \varepsilon v] \ge I[u_*]$ for all $\varepsilon$ and $v \in C_0^2(\Omega)$.}

    \item (2 points) Show that when $\Omega = [a, b]$, $u_*$ must be the straight line connecting $(a, g(a))$ and $(b, g(b))$.
\end{enumerate}
\textbf{[Total for Question 6: 5 points.]}

\noindent \rule[0.5pt]{\textwidth}{0.5pt}

\textbf{Question 7.}
\begin{enumerate}[label=(\alph*)]
    \item (2 points) Suppose $u : \mathbb{R}^2 \to \mathbb{R}^2$ is given by
          \[
              u(r \cos \theta, r \sin \theta) = R(r)\Phi(\theta).
          \]
          Compute $\Delta u$ in polar coordinates $(r, \theta)$.

    \item (3 points) Use the method of Separation of Variables in polar coordinates to derive a solution formula to the Laplace equation
          \[
              \Delta u = 0, \quad x \in B_1(0), \qquad u(x) = g(x), \quad x \in \partial B_1(0),
          \]
          where
          \[
              g(\cos \theta, \sin \theta) = \sum_{n=0} a_n \cos(n\theta) + \sum_{n=1}^{\infty} b_n \sin(n\theta).
          \]
\end{enumerate}
\textbf{[Total for Question 7: 5 points.]}

\end{document}
